This systematic review synthesized 92 studies investigating the neurophysiology of cybersickness in VR.
By uniting neuroimaging with non-invasive neurostimulation under a single neuromarker framework, this review highlights the field's shift away from retrospective subjective questionnaires and toward continuous, objective indices of user state.

The emerging evidence points to specific spectral signatures of cybersickness, primarily characterized by increased low-frequency (delta, theta) power and a suppression of alpha-band oscillations in the electroencephalogram.
However, these responses are entangled with the hardware used to elicit them.
Our meta-analytic findings demonstrate that VR display hardware substantially modulates alpha- and beta-band responses, indicating that the hardware configuration itself is a critical confound that must be accounted for when interpreting neural correlates of cybersickness.

Progress toward generalizable neuromarkers is currently hindered by methodological fragmentation and sparse reporting.
A methodological audit of the classification literature revealed a lack of cross-subject validation, undisclosed decision thresholds, and an over-reliance on accuracy metrics under class imbalance.
These barriers limit the replicability of classification models and impede the synthesis of findings across laboratories.
Addressing these limitations requires a concerted shift toward open science practices, standardized benchmarks, and interventional designs that test the causal relevance of observed neural patterns.

Overcoming these barriers is essential for realizing the full potential of VR.
Establishing robust, generalizable neuromarkers and testing them through closed-loop mitigation strategies will pave the way for adaptive, user-aware immersive systems that can continuously monitor and respond to user discomfort.
This line of research is critical for making immersive experiences universally comfortable and accessible.
